\subsection{Proposed Solution}

To overcome the limitations of high coupling, rigid control flow, and poor maintainability inherent in the naïve conditional approach, we propose restructuring the system using the \textbf{Chain of Responsibility (CoR)} behavioral design pattern \cite{gof1994}. This pattern introduces a more modular and extensible mechanism for processing requests through a sequence of independent components.

\subsubsection{Core Concept and Architectural Idea}
The fundamental principle of the Chain of Responsibility pattern is to decouple the sender of a request from its potential receivers. Rather than embedding all validation logic within a single monolithic structure, the processing workflow is decomposed into a series of independent objects known as \textbf{handlers}, each responsible for a specific task.

These handlers are organized into a sequential chain, forming a processing pipeline. When a request enters the system:

\begin{enumerate}
    \item The request is first received by the initial handler in the chain.
    \item Each handler evaluates the request based on its own responsibility and determines whether to process, transform, or reject it.
    \item If the request remains valid, it is forwarded to the next handler via a reference, continuing the chain execution.
\end{enumerate}

This delegation-based mechanism eliminates the need for explicit control flow logic in the client and promotes a cleaner separation of concerns.

\subsubsection{Architectural Components}
The proposed architecture consists of four primary components:

\begin{itemize}
    \item \textbf{Handler (Interface / Abstract Class):} Defines a common interface for handling requests and for linking to the next handler in the chain.
    
    \item \textbf{BaseHandler (Optional):} Provides a default implementation for chaining behavior by maintaining a reference to the next handler and forwarding the request when appropriate.
    
    \item \textbf{Concrete Handlers:} Specialized subclasses that encapsulate distinct processing logic (e.g., \texttt{\allowbreak AuthenticationHandler}, \texttt{AuthorizationHandler}, \texttt{ValidationHandler}), each focusing on a single responsibility.
    
    \item \textbf{Client:} Responsible for constructing and configuring the handler chain dynamically, and for initiating the request processing by invoking the first handler.
\end{itemize}

This modular decomposition significantly improves extensibility and reusability across different system contexts.

\begin{figure}[H]
    \centering
    \begin{tikzpicture}[
        node distance=1.8cm and 2.0cm,
        classbox/.style={
            draw, 
            rectangle, 
            minimum width=4.2cm, 
            minimum height=1.8cm, 
            align=center, 
            fill=blue!5, 
            font=\small\sffamily
        },
        arrow/.style={-latex, thick},
        inherit/.style={-{Triangle[open, length=3mm, width=3mm]}, thick}
    ]

    % Nodes
    \node (base) [classbox] {
        \textbf{\textit{BaseHandler}}\\[2pt]
        \rule{3.8cm}{0.4pt}\\[3pt]
        \footnotesize-- nextHandler: BaseHandler\\[1pt]
        \footnotesize+ setNext(h): BaseHandler\\[1pt]
        \footnotesize+ handle(r: Request)
    };
    
    \node (client) [classbox, left=1.8cm of base, minimum height=1.3cm] {
        \textbf{Client / Application}\\[2pt]
        \rule{3.8cm}{0.4pt}\\[3pt]
        \footnotesize+ main()
    };
    
    \node (auth) [classbox, below left=1.8cm and 0.2cm of base, minimum height=1.3cm] {
        \textbf{AuthenticationHandler}\\[2pt]
        \rule{3.8cm}{0.4pt}\\[3pt]
        \footnotesize+ handle(r: Request)
    };
    
    \node (valid) [classbox, below right=1.8cm and 0.2cm of base, minimum height=1.3cm] {
        \textbf{ValidationHandler}\\[2pt]
        \rule{3.8cm}{0.4pt}\\[3pt]
        \footnotesize+ handle(r: Request)
    };

    % Connections
    \draw[inherit] (auth) -- (base);
    \draw[inherit] (valid) -- (base);
    \draw[arrow] (client) -- node[above, font=\scriptsize] {initiates} (base);
    
    % FIX MŨI TÊN TỰ TRỎ: Cho đi từ cạnh phải (east) vòng sang cạnh trên (north) với khoảng cách rộng rãi
    \draw[arrow] (base.east) -- ++(1.2, 0) -- ++(0, 1.5) node[above, pos=0.5, font=\tiny] {nextHandler} -| ([xshift=1.2cm]base.north);

    \end{tikzpicture}
    \caption{UML Class Diagram for Chain of Responsibility Pattern (Authors' illustration)}
    \label{fig:cor_class_diagram}
\end{figure}

\subsubsection{Refactored Implementation}
The following C++ implementation demonstrates how the Chain of Responsibility pattern replaces monolithic conditional logic with a flexible, object-oriented structure:

\begin{lstlisting}[language=C++, caption={Chain of Responsibility Pattern Implementation}]
#include <iostream>
#include <memory>

// Request definition
struct Request {
    bool isAuthenticated;
    bool isValid;
};

// 1. Base Handler Abstract Class
class BaseHandler {
private:
    std::shared_ptr<BaseHandler> nextHandler;

public:
    virtual ~BaseHandler() = default;

    std::shared_ptr<BaseHandler> setNext(std::shared_ptr<BaseHandler> next) {
        this->nextHandler = next;
        return next; // Enables fluent chaining
    }

    virtual void handle(const Request& request) {
        if (nextHandler) {
            nextHandler->handle(request);
        }
    }
};

// 2. Concrete Handlers
class AuthenticationHandler : public BaseHandler {
public:
    void handle(const Request& request) override {
        if (!request.isAuthenticated) {
            std::cout << "[FAIL] Authentication failed. Terminating chain.\n";
            return;
        }
        std::cout << "[SUCCESS] Authentication passed.\n";
        BaseHandler::handle(request);
    }
};

class ValidationHandler : public BaseHandler {
public:
    void handle(const Request& request) override {
        if (!request.isValid) {
            std::cout << "[FAIL] Data validation failed. Terminating chain.\n";
            return;
        }
        std::cout << "[SUCCESS] Data validation passed.\n";
        BaseHandler::handle(request);
    }
};

class CoreLogicHandler : public BaseHandler {
public:
    void handle(const Request& request) override {
        // Execute actual core logic
        std::cout << "Executing business logic...\n";
        BaseHandler::handle(request);
    }
};

// 3. Client Configuration and Execution
int main() {
    auto authHandler = std::make_shared<AuthenticationHandler>();
    auto validHandler = std::make_shared<ValidationHandler>();
    auto coreHandler = std::make_shared<CoreLogicHandler>();

    // Build the chain
    authHandler->setNext(validHandler)->setNext(coreHandler);

    Request userRequest = {true, true};
    authHandler->handle(userRequest);

    return 0;
}
\end{lstlisting}

\begin{figure}[H]
    \centering
    \begin{tikzpicture}[
        lifeline/.style={draw, dashed, gray!80, thick},
        entity/.style={
            draw, 
            rectangle, 
            fill=blue!10, 
            minimum width=2.6cm, 
            minimum height=0.8cm, 
            font=\small\sffamily\bfseries,
            align=center
        },
        msg/.style={-latex, thick, font=\xsf},
        ret/.style={-latex, dashed, thick, font=\xsf}
    ]

    % Define font size shortcut
    \newcommand{\xsf}{\fontsize{8}{10}\sffamily}

    % Lifeline entities
    \node (client) at (0,0) [entity] {:Client};
    \node (auth) at (3.5,0) [entity] {:AuthHandler};
    \node (val) at (7.0,0) [entity] {:ValidationHandler};
    \node (core) at (11.0,0) [entity] {:CoreLogicHandler};

    % Draw Lifelines
    \draw [lifeline] (client) -- (0,-6.5);
    \draw [lifeline] (auth) -- (3.5,-6.5);
    \draw [lifeline] (val) -- (7.0,-6.5);
    \draw [lifeline] (core) -- (11.0,-6.5);

    % Messages and Execution Flow
    % 1. Client calls AuthHandler
    \draw [msg] (0,-0.9) -- node[above, font=\xsf] {handle(request)} (3.5,-0.9);
    \node[anchor=west, font=\tiny, color=blue!80!black] at (3.6, -1.3) {[Authentication Passed]};
    
    % 2. AuthHandler delegates to ValidationHandler
    \draw [msg] (3.5,-2.0) -- node[above, font=\xsf] {BaseHandler::handle(request)} (7.0,-2.0);
    \node[anchor=west, font=\tiny, color=blue!80!black] at (7.1, -2.4) {[Validation Passed]};
    
    % 3. ValidationHandler delegates to CoreLogicHandler
    \draw [msg] (7.0,-3.1) -- node[above, font=\xsf] {BaseHandler::handle(request)} (11.0,-3.1);
    
    % 4. CoreLogicHandler executes logic
    \draw [msg] (11.0,-4.0) -- ++(1.4,0) -- ++(0,-0.6) -- node[below, pos=0.85, font=\tiny] {Execute Business Logic} (11.0,-4.6);

    % 5. Return calls
    \draw [ret] (11.0,-5.2) -- (7.0,-5.2);
    \draw [ret] (7.0,-5.6) -- (3.5,-5.6);
    \draw [ret] (3.5,-6.0) -- (0,-6.0);

    \end{tikzpicture}
    \caption{UML Sequence Diagram showing successful request execution flow (Authors' illustration)}
    \label{fig:cor_sequence_diagram}
\end{figure}

\subsubsection{Resolution of Design Limitations}
By adopting the Chain of Responsibility architecture, the system effectively resolves the structural shortcomings of the traditional approach:

\begin{itemize}
    \item \textbf{Single Responsibility Principle (SRP) \cite{martin2000design}:} Each handler encapsulates a single processing concern, leading to clearer separation of logic and improved maintainability.
    
    \item \textbf{Open/Closed Principle (OCP) \cite{martin2000design}:} New processing steps can be introduced by extending the handler hierarchy without modifying existing, stable components.
    
    \item \textbf{Loose Coupling:} The client interacts only with the head of the chain, remaining independent of the internal composition and implementation details of subsequent handlers.
    
    \item \textbf{Runtime Flexibility:} The processing chain can be dynamically assembled, reordered, or extended at runtime, enabling adaptive system behavior based on context or configuration.
\end{itemize}

Overall, this design transforms a rigid and tightly coupled structure into a flexible, scalable, and maintainable processing pipeline.

The following section further evaluates the effectiveness of this pattern through practical analysis and real-world applications.